\section{Theory: Role-Label-Conditioned Routing}
\label{sec:theory}

\subsection{The Routing Problem in Heterogeneous MARL}

In cooperative MARL with heterogeneous roles, agents must not only follow their own behavioral mandate but also identify \emph{which teammates} are most relevant to their decision at each timestep. This is a routing problem: given a set of ally entities, which subset should a given agent condition on, and how strongly?

Role labels provide a natural signal for this selection. A ranged attacker (Strike) benefits most from knowing the position of the frontline tank (Vanguard), not the healer. A healer (Support) prioritizes the most at-risk ally. These routing preferences are determined by role-pair relationships, not spatial proximity alone. The central question of this paper is: \emph{do different policy architectures differ in how effectively they exploit role labels as routing signals?}

\subsection{MLP Cannot Structurally Encode Role-to-Role Routing}

In an MLP policy, the role label is concatenated with all other observation features and processed by shared weight matrices:
\[
h = \text{MLP}([o_{\text{pos}},\ o_{\text{hp}},\ r_{\text{onehot}},\ \ldots])
\]
The computation ``which entity to condition on given my role'' is not structurally encoded. It must be implicitly learned as a feature interaction buried in weight products $W_1 W_2 \cdots W_n$. The MLP has no structural inductive bias toward routing by role-pair. Gradient attribution confirms that role-label signal is much weaker than spatial signal in the MLP baseline (Section~\ref{sec:experiments}), while the attention mechanism produces markedly more structured routing patterns.

Note that this does not claim MLP routing is structurally impossible. The claim is calibrated: attention provides a \emph{more direct structural inductive bias} for role-pair routing. Whether this bias activates is verified empirically.

\subsection{Why Self-Attention Provides a Direct Routing Path}

In the proposed architecture, role labels appear in two positions in the observation:
\begin{itemize}
  \item \texttt{self\_tok[4:7]}: the agent's own role one-hot, encoded into the Query.
  \item \texttt{ally\_tok[12:15]}, \texttt{ally\_tok[20:23]}: each ally's role one-hot, encoded into the corresponding Key/Value.
\end{itemize}

The cross-attention score between agent $i$ (role $r_i$) and ally $j$ (role $r_j$) is:
\[
\text{attn}_{ij} = \text{softmax}\!\left(\frac{(W_Q \cdot \mathbf{x}_i)(W_K \cdot \mathbf{x}_j)^\top}{\sqrt{d_k}}\right)
\]
where $\mathbf{x}_i$ includes $r_i$ and $\mathbf{x}_j$ includes $r_j$. Because role labels condition both Q and K simultaneously, the attention score is a \emph{direct function of the role-pair $(r_i, r_j)$}, not just spatial features. The network can learn role-to-role routing via weight matrices $W_Q$ and $W_K$ that project the role one-hot into the query and key spaces respectively.

This computation is short (one matrix multiplication per head) and attributable: the routing can be read out directly from attention weights, as done in Section~\ref{sec:experiments}.

\subsection{Intra-Set Self-Attention: Relational Context Before Routing}

Before cross-attention, each entity group passes through an Intra-Set Self-Attention block. This provides relational context within the ally set before the agent's query is formed. Specifically, the ally token for Vanguard can be enriched by its relationship to the Support token before entering the cross-attention Key computation.

This explains the performance gap between Condition 1 (Full) and Condition 2 (No Self-Attn): without intra-set contextualization, Keys and Values entering cross-attention are computed from raw entity embeddings without relational structure. Intra-set self-attention contributes beyond the role-routing path, but the routing mechanism itself is primarily located in cross-attention.

\subsection{Testable Predictions}

The structural argument generates four testable predictions, evaluated in Section~\ref{sec:experiments}:

\begin{enumerate}
  \item Attention produces more structured role routing than MLP (routing matrix std is higher). \textbf{[PASS]}
  \item Removing role labels degrades performance by $>$5pp when a shared policy is used. \textbf{[PASS]}
  \item Formation sensitivity ($\Delta_{\max}$) is highest for Vanguard (frontline, formation-sensitive). \textbf{[REVISED: Support highest]}
  \item Strike attends preferentially to Support (healer-seeking behavior). \textbf{[REVISED: Strike$\to$Vanguard dominant]}
\end{enumerate}

Predictions 3 and 4 are revised: the emergent routing pattern differs from the domain-knowledge prediction but remains tactically coherent. This is the paper's central translation-gap result rather than a failure case.
